%=====================================================================
\chapter{The Relational Model}\label{ch:relational}
%=====================================================================
\locator{2}{Part II --- Designing It}

\begin{outcomes}
By the end of this chapter you will be able to:
\begin{tight}
  \item \textbf{Define} relation, tuple, attribute, domain, degree and
        cardinality, and \textbf{use} each word correctly.
  \item \textbf{State} the four properties of a relation and \textbf{show}
        which of them a SQL table does not have.
  \item \textbf{Distinguish} a relation schema from a relation instance, and
        \textbf{write} both in the standard notation.
  \item \textbf{Identify} superkeys, candidate keys and the primary key of a
        relation from its data \emph{and} its guarantees.
  \item \textbf{Explain} what a null value is, what it is not, and why it is
        the model's weakest point.
\end{tight}
\end{outcomes}

\begin{prereq}
\chref{ermodel} for keys --- the definitions are the same ones, restated for
relations. \chref{firstdatabase} for the \relation{}{} notation and for
\texttt{sqlite3}, which you will use at the end to break one of this
chapter's rules on purpose.
\end{prereq}

\begin{keyterms}
relational data model $\bullet$ relation $\bullet$ tuple $\bullet$ attribute
$\bullet$ domain $\bullet$ atomic value $\bullet$ relation schema $\bullet$
relation instance $\bullet$ degree $\bullet$ arity $\bullet$ cardinality
$\bullet$ relational database schema $\bullet$ superkey $\bullet$ candidate
key $\bullet$ primary key $\bullet$ foreign key $\bullet$ null value
$\bullet$ unknown $\bullet$ inapplicable
\end{keyterms}

\section{One Idea, and It Is Mathematics}

Codd's 1970 paper proposed that all data in a database be represented one way:
as relations. Not as files, not as linked records, not as trees. One structure,
with a mathematical definition, so that questions about data could be answered
by reasoning rather than by knowing how the file was laid out.

That is the whole of the relational model, and the reason every system in
\chref{whynotfiles}'s list --- Oracle, PostgreSQL, MySQL, SQL Server, SQLite
--- is a different implementation of the same five pages of theory.

\begin{definitionbox}[Domain, attribute, tuple, relation]
A \term{domain} is a named set of \term{atomic} values --- values the model
treats as indivisible. \texttt{semester} might have domain
$\{1,2,\dots,8\}$; \texttt{roll\_no} the set of strings matching the
university's pattern.

\smallskip
An \term{attribute} is a name for the role a domain plays in a relation. Two
attributes may share a domain: \texttt{start\_year} and \texttt{end\_year} are
both years and are not the same attribute.

\smallskip
A \term{tuple} is an ordered list of values, one from each attribute's domain.
Informally, a row.

\smallskip
A \term{relation} is a \emph{set} of tuples over the same attributes.
Informally, a table --- and the gap between ``informally'' and ``exactly'' is
section~7.3.
\end{definitionbox}

\begin{examplebox}[The vocabulary on one small relation]
\relation{student}{\relkey{roll\_no}, name, programme, semester}

\medskip
\begin{tabular}{@{}L{3.2cm}L{3.0cm}L{2.4cm}L{2.2cm}@{}}
\textbf{roll\_no} & \textbf{name} & \textbf{programme} & \textbf{semester} \\
2021-IT-001 & Ayesha Khan & BSIT & 2 \\
2021-IT-002 & Bilal Ahmed & BSIT & 2 \\
2021-CS-007 & Hira Nadeem & BSCS & 4 \\
\end{tabular}

\medskip
\begin{tight}
  \item The \textbf{relation schema} is the heading: the name plus the four
        attributes. It is written \relation{student}{\relkey{roll\_no}, name,
        programme, semester}.
  \item The \textbf{relation instance} is the three tuples. Tomorrow it is
        different; the schema is not.
  \item The \textbf{degree} (or \term{arity}) is \textbf{4} --- the number of
        attributes. It changes only when the schema changes.
  \item The \textbf{cardinality} is \textbf{3} --- the number of tuples. It
        changes every time somebody enrols.
  \item The \textbf{domain} of \texttt{semester} is the integers 1 to 8. The
        value 2 appears twice and that is fine; domains are not keys.
\end{tight}

\smallskip
\textbf{Degree and cardinality are easy to swap by accident.} Degree counts
columns and is a property of the schema; cardinality counts rows and is a
property of the instance. ``A relation of degree 4 and cardinality 3'' names
the table above exactly.
\end{examplebox}

\dsfig{%
\begin{tikzpicture}[
  hd/.style={font=\scriptsize\bfseries, text=white},
  cell/.style={font=\scriptsize\ttfamily, text=primary!65!black},
  note/.style={font=\scriptsize\itshape, text=neutral, align=left},
  tag/.style={font=\scriptsize\bfseries, text=accent},
  br/.style={draw=accent!70, line width=0.8pt,
             decorate, decoration={brace, amplitude=4pt}}
]
  % heading band
  \fill[primary, rounded corners=2pt] (0,0) rectangle (9.2,0.55);
  \node[hd, anchor=west] at (0.15,0.28)  {roll\_no};
  \node[hd, anchor=west] at (2.55,0.28)  {name};
  \node[hd, anchor=west] at (5.15,0.28)  {programme};
  \node[hd, anchor=west] at (7.65,0.28)  {semester};
  % body
  \foreach \i/\a/\b/\c/\d in {%
      1/2021-IT-001/Ayesha Khan/BSIT/2,
      2/2021-IT-002/Bilal Ahmed/BSIT/2,
      3/2021-CS-007/Hira Nadeem/BSCS/4}{
    \pgfmathsetmacro{\yy}{-0.42*\i}
    \ifodd\i\fill[primary!5] (0,\yy-0.17) rectangle (9.2,\yy+0.23);\fi
    \node[cell, anchor=west] at (0.15,\yy) {\a};
    \node[cell, anchor=west] at (2.55,\yy) {\b};
    \node[cell, anchor=west] at (5.15,\yy) {\c};
    \node[cell, anchor=west] at (7.65,\yy) {\d};}
  \draw[primary!35, line width=0.8pt] (0,-1.43) rectangle (9.2,0.55);

  % annotations
  \draw[br] (9.45,0.5) -- node[tag, right=5pt] {} (9.45,-1.4);
  \node[tag, anchor=west, align=left] at (9.9,-0.45)
       {cardinality\\= 3 tuples};
  \draw[br] (0,0.85) -- (9.2,0.85);
  \node[tag, anchor=south] at (4.6,1.0) {degree = 4 attributes};

  \node[note, anchor=east, text width=3.3cm] at (-0.35,0.28)
       {\textbf{\color{primary}heading}\\the relation \emph{schema}};
  \node[note, anchor=east, text width=3.3cm] at (-0.35,-0.85)
       {\textbf{\color{primary}body}\\the relation \emph{instance}};
  \node[note, anchor=north west, text width=5.4cm] at (0,-1.75)
       {\textbf{one tuple} --- a value from each attribute's domain, every one
        of them \emph{atomic}};
  \node[note, anchor=north west, text width=5.2cm] at (5.4,-1.75)
       {\textbf{a set} --- no tuple repeats, and no tuple is first};
  \draw[-{Stealth[length=1.8mm]}, neutral!60] (1.6,-1.72) -- (1.3,-1.3);
  \draw[-{Stealth[length=1.8mm]}, neutral!60] (7.4,-1.72) -- (7.4,-1.45);
\end{tikzpicture}}%
{One relation, with every word of the vocabulary attached to the thing it
names. Degree belongs to the heading and changes only when the schema does;
cardinality belongs to the body and changes all day.}{relanatomy}

\section{The Four Properties of a Relation}

\begin{definitionbox}[What makes a relation a relation]
\begin{tightnum}
  \item \textbf{No duplicate tuples.} A relation is a \emph{set}. The same
        tuple cannot appear twice --- which is another way of saying every
        relation has a key, if only the whole tuple.
  \item \textbf{Tuples are unordered.} There is no first row. Two relations
        with the same tuples in a different sequence are the same relation.
  \item \textbf{Attributes are unordered.} Attributes are referred to by name,
        not position. \texttt{(name, roll\_no)} and \texttt{(roll\_no, name)}
        are the same heading.
  \item \textbf{Every value is atomic.} One value per attribute per tuple: no
        lists, no nested tables, no ``three phone numbers in one field''.
\end{tightnum}
\end{definitionbox}

\begin{conceptbox}[Why each property is worth having]
These are not arbitrary tidiness. Each one buys something specific.

\begin{tight}
  \item \textbf{No duplicates} means every fact is stored once and can be
        addressed. Two identical rows cannot be told apart, so neither can be
        updated or deleted individually --- which is the file-based system's
        redundancy problem (\chref{whynotfiles}) reappearing inside the
        database.
  \item \textbf{Unordered tuples} is what lets the optimiser
        (\chref{whodoeswhat}) choose any plan it likes. If row order were part
        of the meaning, no index could ever be used and \chref{indexes} would
        not exist.
  \item \textbf{Unordered attributes} is what lets you add a column without
        breaking programs --- the logical data independence of
        \chref{architecture}, and the reason \sqlkw{SELECT *} is a bad habit:
        it depends on an order the model says is not there.
  \item \textbf{Atomic values} is what makes the query language finite. If a
        field could hold a list, every query would need loops, and relational
        algebra (\chref{algebra}) could not be written down. This property is
        also \textbf{first normal form}, arrived at from the other direction
        in \chref{normalforms}.
\end{tight}
\end{conceptbox}

\begin{alertbox}[A SQL table is not a relation, and the difference is measurable]
SQL was designed around the relational model and departs from it in three
places. You should know all three, because each one is a bug waiting to
happen.

\medskip
\begin{tabular}{@{}L{4.0cm}L{4.6cm}L{5.1cm}@{}}
\textbf{Relation} & \textbf{SQL table} & \textbf{What follows} \\
no duplicate tuples & \textbf{duplicates allowed} unless a key is declared &
  a table without a key is not a relation at all \\
tuples unordered & \sqlkw{ORDER BY} imposes one, and without it the order is
  arbitrary & never rely on the order you happen to see \\
values atomic & mostly enforced, but a text column will hold
  \texttt{"0300-1234567, 0321-9876543"} quite happily &
  the model cannot stop you; only you can \\
\end{tabular}

\smallskip
The first is the serious one and you can demonstrate it in four lines:

\begin{lstlisting}[language=sql]
CREATE TABLE t (a TEXT, b TEXT);        -- no key declared
INSERT INTO t VALUES ('x','y'), ('x','y');
SELECT COUNT(*) FROM t;                 -- 2
DELETE FROM t WHERE a='x' AND b='y';    -- deletes BOTH
\end{lstlisting}

Two identical rows, and no \sqlkw{WHERE} clause in the language can address one
of them, because there is nothing that distinguishes them. \textbf{This is why
every table in this course declares a primary key}, and \chref{constraints} is
about making the declaration real.
\end{alertbox}

\section{Keys, Restated for Relations}

The definitions are \chref{ermodel}'s, which is the point: the ER model and the
relational model agree about keys, and that agreement is what makes
\chref{mapping} possible.

\begin{definitionbox}[Keys of a relation]
A \term{superkey} is a set of attributes $K$ such that no two tuples of any
legal instance agree on all of $K$.

\smallskip
A \term{candidate key} is a minimal superkey: no proper subset of it is a
superkey.

\smallskip
The \term{primary key} is the chosen candidate key. Its attributes may not be
null --- an identifier that is missing identifies nothing.

\smallskip
A \term{foreign key} is a set of attributes in one relation whose values must
appear as the primary key of another (or be null). It is how relations are
joined to each other, and it is \chref{constraints}'s subject.
\end{definitionbox}

\begin{worked}[Reading keys off an instance, and why you cannot]
\begin{tabular}{@{}L{2.8cm}L{2.6cm}L{2.2cm}L{2.0cm}L{2.4cm}@{}}
\textbf{roll\_no} & \textbf{name} & \textbf{prog} & \textbf{sem}
& \textbf{email} \\
2021-IT-001 & Ayesha Khan & BSIT & 2 & ayesha@iub.pk \\
2021-IT-002 & Bilal Ahmed & BSIT & 2 & bilal@iub.pk \\
2021-CS-007 & Hira Nadeem & BSCS & 4 & hira@iub.pk \\
\end{tabular}

\medskip
\textbf{Step 1 --- what is unique in this instance?} All of
\texttt{roll\_no}, \texttt{name}, \texttt{email}; also
$\{$\texttt{prog}, \texttt{sem}$\}$ is \emph{not} --- rows 1 and 2 agree.

\medskip
\textbf{Step 2 --- which of those are superkeys?} This is where the exercise
stops being mechanical. A superkey is a constraint on \emph{every legal
instance}, and three rows cannot establish one.
\begin{tight}
  \item \texttt{roll\_no}: the university guarantees it. \textbf{Superkey.}
  \item \texttt{email}: issued one per student. \textbf{Superkey} --- on that
        guarantee, not on these rows.
  \item \texttt{name}: unique \emph{here}. The university guarantees nothing.
        \textbf{Not a superkey}, and the data cannot tell you so.
\end{tight}

\medskip
\textbf{Step 3 --- minimality.} $\{$\texttt{roll\_no}$\}$ and
$\{$\texttt{email}$\}$ are both minimal, so both are candidate keys.
$\{$\texttt{roll\_no}, \texttt{sem}$\}$ is a superkey and not minimal.

\medskip
\textbf{Step 4 --- choose.} \texttt{roll\_no}, on \chref{ermodel}'s four
grounds: unique, always present, never changing, short. \texttt{email} fails
``never changing''.

\medskip
\textbf{The asymmetry worth remembering.} An instance can \emph{disprove} a key
instantly --- one duplicate and it is settled. It can never prove one. So:
check the data to find candidates, then ask a human to confirm each, and
record the answer. The examination version of this question gives you the
guarantees; the real version does not, and you have to go and ask.
\end{worked}

\section{Null, and Why It Is the Awkward Part}

\begin{definitionbox}[Null]
A \term{null} is a marker in a tuple meaning that no value from the domain is
recorded for that attribute. It is \textbf{not} zero, \textbf{not} the empty
string, and \textbf{not} a value of any kind --- it is the absence of one.

\smallskip
Codd distinguished two reasons for a null and most systems, including SQLite,
distinguish neither:
\begin{tight}
  \item \term{unknown} --- the student has a CNIC, we do not know it;
  \item \term{inapplicable} --- the student has no CNIC, because they are
        fifteen.
\end{tight}
One marker, two meanings, and no way to tell them apart in a query.
\end{definitionbox}

\begin{examplebox}[Three nulls in one row, meaning three different things]
\begin{tabular}{@{}L{2.9cm}L{3.4cm}L{7.6cm}@{}}
\textbf{Attribute} & \textbf{Value} & \textbf{What the null means} \\
\texttt{cnic}        & \textsc{null} & unknown --- they have one, it is not on file \\
\texttt{thesis\_title} & \textsc{null} & inapplicable --- second-semester
                                        students do not write one \\
\texttt{grade}       & \textsc{null} & not yet --- the semester is not over \\
\end{tabular}

\smallskip
All three print as blank and all three behave identically in every comparison.
A query that counts ``students with no CNIC on file'' will silently include
anybody whose CNIC was inapplicable, and nothing in the database records that
the two cases were different.

\smallskip
\textbf{What a designer does about it.} Where the distinction matters, make it
explicit: a separate \texttt{cnic\_status} attribute, or a different relation
for thesis students. Where it does not, allow the null and write the ambiguity
down. The one thing not to do is invent a magic value --- \texttt{'N/A'},
\texttt{0000-0000000}, \texttt{9999-12-31} --- which is a null that the
database does not know is a null, and which will be counted, averaged and
sorted as though it were real.
\end{examplebox}

\begin{alertbox}[Null is ahead of you in three later chapters]
This chapter only names the problem. It returns, each time worse:

\begin{tight}
  \item \chref{select}: \texttt{WHERE grade = NULL} matches nothing, ever ---
        not even the rows where grade is null. Comparison with null is neither
        true nor false but \emph{unknown}, and SQL's logic has three values.
  \item \chref{grouping}: \sqlkw{COUNT(*)} and \sqlkw{COUNT(grade)} give
        different answers on the same rows, and \sqlkw{AVG} silently ignores
        nulls rather than treating them as zero --- which is usually right and
        occasionally catastrophic.
  \item \chref{joins}: a row whose foreign key is null disappears from an inner
        join, so the section with no instructor vanishes from the timetable
        report and nobody notices until a Tuesday.
\end{tight}

\smallskip
Nulls are the single largest source of quietly wrong answers in SQL. Every one
you allow in a schema is a decision; allow them deliberately.
\end{alertbox}

\begin{pitfall}
\begin{tight}
  \item \textbf{Saying ``table'' when the question asks about a relation.} In
        an examination the distinction is the answer: a table may have
        duplicates and an order.
  \item \textbf{Swapping degree and cardinality.} Degree counts attributes,
        cardinality counts tuples.
  \item \textbf{Relying on row order.} Without \sqlkw{ORDER BY} there is none,
        and the order you see today comes from the plan, which
        \chref{indexes} will change.
  \item \textbf{Reading keys off sample data.} Data disproves keys; it never
        proves them.
  \item \textbf{Treating null as zero or as the empty string.} It is neither,
        and the arithmetic will show it.
  \item \textbf{Inventing a magic value instead of a null.} It gets averaged.
  \item \textbf{Creating a table with no primary key.} It is not a relation,
        and you have given up the ability to address a row.
\end{tight}
\end{pitfall}

\begin{lab}[Break three properties of a relation, deliberately]
\begin{lstlisting}[language=bash]
# --- 1. A scratch database; keep university.db clean ----------------
sqlite3 relmodel.db
\end{lstlisting}
\begin{lstlisting}[language=sql]
-- --- 2. A table that is not a relation --------------------------
.headers on
.mode box
CREATE TABLE nokey (roll_no TEXT, name TEXT);
INSERT INTO nokey VALUES ('2021-IT-001','Ayesha Khan');
INSERT INTO nokey VALUES ('2021-IT-001','Ayesha Khan');
SELECT COUNT(*) FROM nokey;
-- Property 1 broken. Now try to delete exactly one of them.
DELETE FROM nokey WHERE roll_no = '2021-IT-001' AND name = 'Ayesha Khan';
SELECT COUNT(*) FROM nokey;
-- Write down how many rows went, and say why you could not keep one.

-- --- 3. The same table, as a relation ---------------------------
CREATE TABLE withkey (roll_no TEXT PRIMARY KEY, name TEXT NOT NULL);
INSERT INTO withkey VALUES ('2021-IT-001','Ayesha Khan');
INSERT INTO withkey VALUES ('2021-IT-001','Ayesha Khan');
-- Record the error. One keyword is the whole difference.

-- --- 4. Prove tuples are unordered ------------------------------
INSERT INTO withkey VALUES ('2021-CS-007','Hira Nadeem');
INSERT INTO withkey VALUES ('2021-IT-002','Bilal Ahmed');
SELECT * FROM withkey;
CREATE INDEX idx_nm ON withkey(name);
SELECT * FROM withkey;
-- Same relation, same query. Did the order change? Either way,
-- say what you are entitled to assume about it.

-- --- 5. Break atomicity, which nothing will stop you doing ------
CREATE TABLE badphone (roll_no TEXT PRIMARY KEY, phones TEXT);
INSERT INTO badphone VALUES
  ('2021-IT-001','0300-1234567, 0321-9876543');
-- Now answer: which students have a number starting 0321?
SELECT roll_no FROM badphone WHERE phones LIKE '%0321%';
-- It appears to work. Now explain, in one sentence, why
-- '0321' in the MIDDLE of a number would also match, and what
-- you would have to write to do this correctly.

-- --- 6. Three nulls, one appearance -----------------------------
CREATE TABLE nulls (roll_no TEXT PRIMARY KEY, cnic TEXT, grade TEXT);
INSERT INTO nulls VALUES ('2021-IT-001', NULL, NULL);
SELECT * FROM nulls WHERE cnic IS NULL;
SELECT * FROM nulls WHERE cnic = NULL;
-- Two queries, one of them returns nothing. Which, and WHY?
-- This is the whole of Chapter 14 in two lines; write your
-- explanation down now and check it when you get there.
\end{lstlisting}

\smallskip
\textbf{What to notice.} Step 2 is the one to keep. Two identical rows, and the
language has no way to refer to one of them --- not because SQLite is
deficient, but because there is genuinely nothing that tells them apart. Every
primary key you declare for the rest of this course is buying your way out of
that position.
\end{lab}

\begin{checkpoint}
\begin{tightnum}
  \item A relation has degree 5 and cardinality 200. Say precisely what each
        number counts, which of the two can change without a schema change,
        and what operation changes the other.
  \item Give a specific query that returns a different answer depending on
        whether duplicate tuples are permitted, and say what you would declare
        to make the answer reliable.
  \item \texttt{grade} is null for a student who has not sat the examination
        and also for one whose paper was lost. Explain why the database cannot
        distinguish them, give two different designs that could, and say what
        each design costs.
\end{tightnum}
\end{checkpoint}

\begin{chaptersummary}
\begin{tight}
  \item A \textbf{domain} is a set of atomic values; an \textbf{attribute}
        names a domain's role; a \textbf{tuple} picks one value per attribute;
        a \textbf{relation} is a \emph{set} of tuples.
  \item \textbf{Degree} counts attributes and belongs to the schema;
        \textbf{cardinality} counts tuples and belongs to the instance.
  \item Four properties: \textbf{no duplicate tuples}, \textbf{tuples
        unordered}, \textbf{attributes unordered}, \textbf{values atomic}.
  \item Each property buys something: addressability, the optimiser's freedom,
        logical data independence, and a finite query language.
  \item \textbf{A SQL table is not a relation.} It permits duplicates unless a
        key is declared, exposes an order, and cannot stop you putting a list
        in a text column.
  \item A table with no primary key gives up the ability to address a row: two
        identical rows cannot be told apart by any \sqlkw{WHERE} clause.
  \item \textbf{Superkey}, \textbf{candidate key} and \textbf{primary key} mean
        in the relational model exactly what they meant in the ER model, which
        is what makes \chref{mapping} mechanical.
  \item An instance can \textbf{disprove} a key but never prove one. Use the
        data to find candidates and a human to confirm them.
  \item A \textbf{null} marks the absence of a value. It is not zero and not
        the empty string, it conflates \emph{unknown} with
        \emph{inapplicable}, and it is the source of quietly wrong answers in
        Chapters~\ref{ch:joins}, \ref{ch:select} and \ref{ch:grouping}.
  \item Never substitute a magic value for a null. It will be counted, averaged
        and sorted as though it were real.
\end{tight}
\end{chaptersummary}

\begin{reviewq}
  \item \textbf{(MCQ)} The number of attributes in a relation is its:
        (a)~cardinality; (b)~degree; (c)~domain; (d)~arity of its key.
  \item \textbf{(MCQ)} Which property of a relation does a SQL table most
        clearly fail to have? (a)~attributes unordered; (b)~values atomic;
        (c)~no duplicate tuples; (d)~a named schema.
  \item \textbf{(Short)} Two attributes share the domain ``year''. Explain why
        they are still two different attributes, with an example.
  \item \textbf{(Short)} Why can an instance disprove a candidate key but never
        establish one? Illustrate with \texttt{name}.
  \item \textbf{(Short)} State the difference between \texttt{WHERE cnic IS
        NULL} and \texttt{WHERE cnic = NULL}, and say which returns rows.
  \item \textbf{(Applied)} You are given a table of 50{,}000 rows and told only
        that ``it is the student list''. Describe the queries you would run to
        propose a primary key, say exactly what each result would and would not
        establish, and state what you would have to ask a person before
        declaring one.
  \item \textbf{(Applied)} A system stores a student's phone numbers in one
        comma-separated text column. List four distinct questions the
        university cannot reliably answer in that design, show the relation you
        would replace it with, and explain which of the four properties of a
        relation the original violated.
\end{reviewq}
